%% 02-gas-lift-surrogate-optimization — Parte I (massa de dados e substitutos). Construção: `make` nesta pasta (ver ../README.md).
%% Guia de autores da revista alvo: ver ../00-escolha-da-revista.md (citações autor–ano, resumo <= 250 palavras,
%% 3–5 highlights <= 85 caracteres, 1–7 palavras-chave, SI, figuras vetoriais, declarações e CRediT no fim).
%% Estado (2026-10-19): um só artigo (decisão do dono): modelo fenomenológico em resumo, massa de dados, NN e PINN e resultados
%% escritos; a secção da otimização operacional está em espera (formulação escrita, resultados por fazer, marcada com \tbd).

\documentclass[a4paper,fleqn]{cas-sc}
%\documentclass[a4paper,fleqn,review]{cas-sc}   % revisão: espaçamento largo e números de linha

\usepackage[authoryear,longnamesfirst]{natbib}
\usepackage{amsmath,amssymb}
\usepackage{booktabs}
\usepackage{siunitx}
\usepackage{placeins}

\graphicspath{{figures/generated/}{figures/static/}}

\newcommand{\aG}{\alpha_G}
\newcommand{\btheta}{\boldsymbol{\theta}}
\newcommand{\bx}{\mathbf{x}}
\newcommand{\by}{\mathbf{y}}
\usepackage{xcolor}
\newcommand{\tbd}[1]{\textcolor{red}{#1}}

\begin{document}
\let\WriteBookmarks\relax
\def\floatpagepagefraction{1}
\def\textpagefraction{.001}

\shorttitle{Operating a gas-lifted well under severe slugging with a two-fluid-model surrogate}
\shortauthors{J.N.M. de Souza and A.L.H. Costa}

\title[mode = title]{Operating a gas-lifted well under severe slugging: time-averaged signatures from a two-fluid simulator, physics-informed surrogates of the choke and gas-lift decisions, and their operational optimisation}

\author[1]{Jaime Neiva Miranda de Souza}

\author[2]{Andr\'e Luiz Hemerly Costa}
\cormark[1]
\ead{andrehc@uerj.br}

\affiliation[1]{organization={Miradouro Digital}}  % TODO: morada, cidade e país

\affiliation[2]{organization={Instituto de Qu\'imica, Universidade do Estado do Rio de Janeiro},
                addressline={Rua S\~ao Francisco Xavier 524},
                city={Rio de Janeiro},
                postcode={20550-900},
                state={RJ},
                country={Brazil}}

\cortext[1]{Corresponding author.}

\begin{abstract}
A gas-lifted satellite well with a flowline and a riser can produce in severe slugging, in a pressure cycle or in a steady state depending on two operating decisions, the opening of the topside choke and the pressure of the lift gas at the platform, and the time-averaged production of an unstable operating point is not what a steady-state model predicts for it. This paper builds the data and the surrogate models of those two decisions and uses them to optimise the operation of the well. The well, flowline, riser and gas-lift branch are simulated with the four-equation two-fluid model of a companion study, with the flow pattern identified dynamically and the gas-lift line, annulus and check valve simulated with the same model and controlled only by the platform gas pressure. A sampling campaign over the two decisions, in rounds of Sobol' points run as warm-started chains along the shortest path from the base case, with an adaptive horizon that detects the established regime and captures three cycles, produced 156 runs in a day on a laptop; each run is reduced to time-averaged markers (mean, minimum, maximum and the two extreme frequencies) of pressure, void fraction and phase flow rates at eight characteristic points of the network. The map of the decisions shows a stability boundary that runs from \SI{80}{\bar} of lift gas with the choke nearly closed to \SI{125}{\bar} with it open, continuous injection modulated by the cycle below it, and a production gain of up to \SI{34}{\percent} above it. Two surrogates of the mean liquid, gas and lift-gas rates and of their spans were trained on the markers: a neural network and a physics-informed network whose physics are the time-averaged mass balances of the well, the gas-lift branch and the valve, exact for the averages and evaluated at collocation points without simulation. The physics raise the test coefficient of determination of the mean liquid rate from 0.57 to 0.88 and those of the gas and lift-gas rates to 0.997, halve the error of the gas rates with ten to forty training runs, predict the maxima of the liquid and lift-gas rates with 0.91 and 0.99, and leave the spans, the signature of slugging, as the quantity that neither network predicts well. The surrogate then enters an operational optimisation of the two decisions: the revenue of the mean liquid rate minus the cost of compressing the lift gas from \SI{5}{\bar} to the platform pressure, with the predicted maxima of the lift-gas and liquid rates bounded by the compressor and separator capacities. At the prices of the example the compression costs less than one per cent of the revenue and the optimum is the corner of maximum gas and open choke, which the simulator confirms to \SI{2}{\percent}; a compressor capacity moves the optimum along the valve law, and the trade-off against the suction pressure \tbd{is left for the next step}.
\end{abstract}

\begin{highlights}
\item Gas-lift branch simulated with the two-fluid model, controlled by the platform gas pressure
\item Warm-started Sobol' campaign with adaptive horizon: 156 runs over choke and lift gas in a day
\item Time-averaged markers at eight points replace trajectories as the training data
\item Physics of the averages lifts the surrogate of the mean liquid rate from $R^2$ 0.57 to 0.88
\item Optimum of revenue minus compression on the surrogate, with capacity constraints on predicted maxima
\end{highlights}

\begin{keywords}
severe slugging \sep gas lift \sep two-fluid model \sep surrogate models \sep physics-informed neural networks \sep design of experiments
\end{keywords}

\maketitle

\section{Introduction}\label{sec:intro}

Continuous gas lift is the artificial-lift method of most deepwater satellite wells in Brazil, and its optimisation --- how much gas to inject in each well of a field that shares a compressor --- has a large literature built on steady-state performance curves \citep{souza2010}. A well with a flowline and a riser, however, does not always have a steady state to be optimised: with the choke open and little gas it produces in severe slugging \citep{schmidt1980,taitel1986,jansen1996}, and the time-averaged production, the liquid surges that reach the separator and the pressure swings at the wellhead are quantities of the cycle, not of the unstable steady state that a nodal model would return for the same point. A companion study \citep{souza2026a} developed a four-equation two-fluid model with dynamic flow-pattern identification that reproduces the cycle in a well--flowline--riser system of field dimensions, located the onset of slugging as the choke opens and showed, with a gas-lift branch added to the system, that injection at 100 to \SI{120}{\bar} regularises the cycle and from \SI{130}{\bar} suppresses it.

This paper is about the operation of that gas-lifted well. Its question is practical: given the two decisions an operator has, the opening of the topside choke and the pressure of the lift gas at the platform, what does the well deliver on average, with what surges and in which regime --- and can a model cheap enough for optimisation answer that over the whole plane of the two decisions? The dynamic simulation is too expensive to be called inside an optimiser (minutes per point), so the paper builds a surrogate, and it does so with a method that follows the companion work on compositional steady-state surrogates \citep{carvalho2026b}: a multilayer perceptron trained on simulated data, and the same network with a physics term \citep{raissi2019,karniadakis2021}. The novelty is on both sides of the surrogate. On the data side, the targets are not profiles but time-averaged signatures of a dynamic simulation, and producing them in affordable time required a campaign design of its own: an adaptive horizon that recognises the established regime, warm starts along a shortest path through the sampled points, and a storage that keeps only what the training needs. On the physics side, the governing equations cannot be imposed on a time average of a cycling system; what can, exactly, are the time-averaged mass balances of the well, of the gas-lift branch and of the valve, and the paper shows that this much physics is what the data alone cannot supply.

Section~\ref{sec:model} summarises the phenomenological model and the gas-lift branch. Section~\ref{sec:data} describes the campaign, the horizon, the markers and the resulting map of the decisions. Section~\ref{sec:surrogates} presents the two surrogates and Section~\ref{sec:results} compares them. Section~\ref{sec:optimisation} formulates the operational optimisation on the surrogate, with capacity constraints on the predicted maxima, and reports its first results.

\section{The phenomenological model}\label{sec:model}

The system is the well--flowline--riser of the companion study \citep{souza2026a} with its gas-lift branch (Fig.~\ref{fig:system}): a vertical well of \SI{1500}{\metre} (\SI{127}{\milli\metre}) fed by a linear inflow performance relation, a flowline of \SI{2000}{\metre} (\SI{152.4}{\milli\metre}) whose last \SI{500}{\metre} descend at \SI{2}{\degree} to a low point, a bend of \SI{5}{\metre} radius, a riser of \SI{317.5}{\metre} and \SI{10}{\metre} of horizontal pipe to a choke that discharges into a separator at \SI{10}{\bar}. The gas-lift branch is a 4-in line that runs from the platform along the route of the production line to the wellhead, the casing--tubing annulus (9~5/8~in by 5~1/2~in) to a gas-lift valve \SI{30}{\metre} above the bottom of the well, and a check valve fed by a chamber at the base of the annulus. The base well has a reservoir pressure of \SI{200}{\bar}, a productivity index of \SI{5}{\cubic\metre\per\day\per\bar} and a gas--oil ratio of \SI{50}{\cubic\metre\per\cubic\metre}; the fluid is the light oil of \citet{carvalho2022} frozen as a polytropic pair.

Every pipe, the gas pipes included, is described by the one-dimensional four-equation two-fluid model of the companion study: mass and momentum of each phase, closures by flow pattern (stratified, intermittent, bubbly, annular) selected dynamically on the local state through smooth logistic switches, a low-void rule, an objective virtual-mass term and an interfacial pressure in the slug unit that keep the model hyperbolic along the simulated trajectories. The staggered finite-volume discretisation with phase-wise donor cells and the conservative variable-step implicit integration of that study are kept unchanged, with a compiled right-hand side and exact Jacobians; the base mesh has 15 cells in the well, 35 in the flowline--riser pipe, 19 in the gas-lift line and 10 in the annulus, 323 states in all. The gas-lift branch is controlled by one number, the gas pressure $P_{gas}$ at the platform end of the line: the injected rate is not prescribed but follows from $P_{gas}$, the static gas columns of the line and the annulus, the valve law $q = 0.0075\,C_V\sqrt{\rho_G\,\Delta P}$ ($C_V = 1.5$, $\Delta P$ in bar, regularised at the origin) and the pressure of the well at the injection point, so that it responds to whatever the well does; at \SI{20}{\bar} the valve never opens and the system is the base case of the companion study with the branch inert. The reservoir gas is the gas of the well; the lift gas adds to it at the injection point.

\begin{figure}[htbp]
\centering
\includegraphics[width=\linewidth]{system}
\caption{The well--flowline--riser system with the gas-lift branch (schematic, from the companion study): the gas-lift line follows the route of the production line, the annulus runs beside the well and the check valve injects \SI{30}{\metre} above the bottom. The two operating decisions are the choke opening $a$ and the platform gas pressure $P_{gas}$.}\label{fig:system}
\end{figure}

\section{Data: campaign, horizon and markers}\label{sec:data}

\subsection{Sampling the two decisions}\label{sec:sampling}

The inputs are the choke opening $a \in [0.1, 1]$ and the platform gas pressure $P_{gas} \in [20, 160]$~bar, with the well fixed at its base state. Points are drawn from a scrambled Sobol' sequence \citep{sobol1967,owen1998} in rounds of 25, so that the rounds together remain a single space-filling design, and each round is run as one chain: the first point is the one closest to the base case (choke open, \SI{20}{\bar}) and the run starts from the final state of the base run of the companion study (\SI{8000}{\second} of severe slugging); every later point starts from the final state of the previous one, in the order of the shortest open path through the 25 points (nearest neighbour improved by 2-opt, on the inputs rescaled to $[-1, 1]$). The warm start spares the transient of a blocked start --- three to five cycles in the companion study --- and keeps the steps between consecutive runs small; a large drop of $P_{gas}$ between neighbours is the one expensive transition, since the gas line and the annulus then drain towards the platform, and the shortest path avoids most of them. Each point is identified by a hash of its input tuple, and only what the training needs is stored: the markers below and the final state that seeds the next run.

\subsection{Adaptive horizon with capture of cycles}\label{sec:horizon}

Each run integrates in windows of \SI{3000}{\second} and monitors two signals, the pressure at the wellhead and the liquid rate through the choke. After a minimum time (\SI{4000}{\second} from a blocked start, \SI{1500}{\second} from a warm start) the regime is taken as established when both signals are permanent (amplitude below \SI{0.1}{\bar} and \SI{0.05}{\kilogram\per\second}) or when two consecutive windows have the same envelope, that is, the same amplitude, mean, minimum and maximum within \SI{10}{\percent} of the amplitude; the envelope criterion, rather than equality of periods, is what accepts the cycles whose blowouts alternate in size. The run then continues for three more cycles of the longest monitored period (\SI{600}{\second} for a steady state), which are the capture window of the markers, and stops; a run that never settles stops at \SI{24000}{\second} and is kept with the markers of its last \SI{4000}{\second} and a flag. A wall-clock budget of \SI{40}{\minute} per run bounds the few points where the two horizontal cells between the riser and the choke switch flow pattern at every step, the one open point of the closures that the companion study reported; those runs are flagged and excluded from the training.

\subsection{Markers at characteristic points}\label{sec:markers}

A run is reduced to markers of its capture window at eight points of the network: the bottom of the well, the injection point, the wellhead, the low point of the flowline, the top of the riser, the two ends of the gas-lift line and the base of the annulus. At each point four signals are recorded --- pressure, void fraction, liquid and gas mass flow rates --- and each signal is summarised by its mean, minimum, maximum and two frequencies, the inverse of the longest and of the shortest interval between successive maxima of the smoothed signal: equal for a periodic cycle, different when blowouts alternate, zero for a permanent signal. The gas injected through the valve is recorded in the same way, and a class is assigned from the wellhead pressure and the choke liquid rate as in the companion study (steady, pressure cycle or severe slugging when the peak liquid rate exceeds twice its mean). Trajectories are not kept: they would cost \SI{12}{\mega\byte} per run against \SI{10}{\kilo\byte} of markers, and the surrogate consumes only the markers.

\subsection{The map of the decisions}\label{sec:map}

Table~\ref{tab:dataset} summarises the campaign and Fig.~\ref{fig:dataset_map} shows it on the plane of the decisions. Of the 156 runs, 83 produce in severe slugging, 10 in a pressure cycle and 63 in a steady state; 88 settled on a repeated cycle, 60 on a steady state, 2 reached the time limit and 6 the budget. The stability boundary is a diagonal that runs from about \SI{80}{\bar} of lift gas with the choke at 0.1 to about \SI{125}{\bar} with it open: the more open the choke, the more gas is needed to suppress the cycle. The check valve opens between 86 and \SI{95}{\bar} almost independently of the choke, so that below \SI{90}{\bar} the data are the base well with the choke as the only decision; between the opening of the valve and the boundary the injection is continuous and modulated by the cycle, which shortens and regularises; above the boundary the well is steady and its production rises with the gas, to \SI{5.3}{\kilogram\per\second} at \SI{160}{\bar} with the choke open, \SI{34}{\percent} above the base case. The peak liquid rate through the choke reaches \SI{45}{\kilogram\per\second} in the slugging region with the choke open and falls below \SI{5}{\kilogram\per\second} in the steady region; the six runs that exhausted the budget sit on the boundary at high opening, where the riser receives the most gas.

\begin{table}[htbp]
\caption{The sampling campaign over the two decisions.}\label{tab:dataset}
\footnotesize
\input{tables/generated/tab_dataset}
\input{tables/generated/tab_dataset_note}
\end{table}

\begin{figure}[htbp]
\centering
\includegraphics[width=\linewidth]{dataset_map}
\caption{The 156 runs on the plane of the choke opening and the platform gas pressure: (a) class of the established regime, with the six runs that exhausted the 40-min integration budget before settling (crosses; their markers come from a partial trajectory and they are excluded from the training); (b) mean liquid rate through the choke; (c) peak liquid rate through the choke on a logarithmic scale.}\label{fig:dataset_map}
\end{figure}

\section{Surrogates}\label{sec:surrogates}

\subsection{Outputs and network}\label{sec:network}

The surrogate maps the two decisions to eleven quantities of the capture window: the mean liquid and gas rates at the top of the riser, the mean gas-lift rate through the valve, the spans (maximum minus minimum) of the same three rates, which are the signature of slugging, the maxima of the liquid rate at the top and of the gas-lift rate, which the capacity constraints of the optimisation use, and three auxiliary mean pressures --- at the bottom of the well, at the injection point and at the base of the annulus --- that the physics of the next section needs. Spans and maxima are learnt as $\log(1 + x)$, so that they are never negative and their orders of magnitude are resolved. The network is a multilayer perceptron \citep{hornik1989} with two hidden layers of 32 hyperbolic-tangent units and a linear output, inputs and outputs standardised with the training statistics, trained with Adam \citep{kingma2015} in full batch for 4000 epochs on the mean squared error of the standardised outputs, with the epochs unrolled in a compiled scan so that a network trains in about a second; 70\,\% of the 150 retained runs train it and 30\,\% test it. A grid of one or two layers, 8 to 64 units and three learning rates (24 networks of each kind, Section~\ref{sec:results}) checked that the architecture is not critical.

\subsection{Physics of the averages}\label{sec:physics}

A physics-informed network adds to the data loss the residual of the governing equations at collocation points that need not be simulated \citep{raissi2019}. Here the targets are time averages of a dynamic system, and the time average of a cycle does not satisfy the two-fluid equations, which are nonlinear; what it satisfies exactly are the balances that are linear in the averaged quantities. Over a periodic cycle or a steady state, the mass of each phase in the system does not change on average, so that
\begin{align}
R_1 &= \bar q_L - J_L\,(P_{res} - \bar P_{bh}), \label{eq:r1}\\
R_2 &= \bar q_G - J_G\,(P_{res} - \bar P_{bh}) - \bar q_{gl}, \label{eq:r2}\\
R_3 &= \bar P_{an} - \kappa\,P_{gas}, \label{eq:r3}\\
R_4 &= \bar q_{gl} - 0.0075\,C_V\sqrt{\rho_G(\bar P)\,\langle \bar P_{an} - \bar P_{inj}\rangle_+}, \label{eq:r4}
\end{align}
where $R_1$ and $R_2$ are the liquid and gas balances between the linear inflow relation at the mean bottom-hole pressure and the mean rates at the top, $R_3$ is the static gas column from the platform to the base of the annulus ($\kappa = 1.138$ for this fluid and depth), and $R_4$ is the valve law written with the mean pressures, the only approximate relation of the four (the mean of a square root is not the square root of the mean, but the modulation of the valve pressure drop is within \SI{25}{\percent} and the error is a few per cent), with a smooth gate that closes the valve when the mean annulus pressure falls below the mean injection pressure. The loss is the data term plus $\lambda$ times the mean square of the four residuals, each scaled by the standard deviation of the quantity it compares, evaluated on 200 Latin-hypercube points \citep{mckay1979} of the plane of the decisions; $\lambda = 1$, and values from 0.1 to 100 give the same fit with residuals decreasing with $\lambda$. The residuals are closed-form functions of the network outputs, so that the gradient of the physics term is exact and cheap, unlike the detached-Jacobian construction that the compositional surrogates needed \citep{carvalho2026b}. The information in the physics is therefore not the dynamics but the accounting: the network is told that whatever enters the well on average leaves it on average, and that the lift gas is set by two mean pressures it also predicts.

\section{Results}\label{sec:results}

Table~\ref{tab:surrogate} and Fig.~\ref{fig:parity} compare the two networks on the 45 test runs. The data-only network predicts the mean gas rate at the top and the mean gas-lift rate with $R^2$ of 0.99 and the three mean pressures with 0.97 to 0.999, but the mean liquid rate only with 0.57: that rate varies little over the plane (3.5 to \SI{5.3}{\kilogram\per\second}) and its time average over three cycles of a surging signal is the noisiest of the targets. The physics-informed network raises it to 0.88, with the root-mean-square error falling from 0.36 to \SI{0.19}{\kilogram\per\second}, because the balance $R_1$ ties it to the mean bottom-hole pressure, which is smooth; the gas and lift-gas rates reach 0.997 and the scaled residuals fall by a factor of two to five. The maxima behave like the means they ride on: the maximum lift-gas rate reaches 0.99 with both networks, since the injection is at most modulated by a quarter, and the maximum liquid rate, which carries the surge, 0.83 and 0.91. The spans are the weak point of both networks: the liquid span reaches 0.75 and 0.88; the gas span and the span of the injected gas stay below 0.5, since the amplitude of a surge in three cycles is a noisy marker by nature. Fig.~\ref{fig:maps} shows the physics-informed surrogate on the plane of the decisions: the mean liquid rate rises with the lift gas above the boundary and falls as the choke closes, the lift-gas rate is zero below \SI{90}{\bar} and rises with $P_{gas}$ as the valve law dictates, and the span of the liquid rate marks the slugging region with a narrow band of artefacts along the boundary, where the data change class within a few bar.

\begin{table}[htbp]
\caption{Test metrics of the two surrogates.}\label{tab:surrogate}
\footnotesize
\input{tables/generated/tab_surrogate}
\input{tables/generated/tab_surrogate_note}
\end{table}

\begin{figure}[htbp]
\centering
\includegraphics[width=\linewidth]{surrogate_parity}
\caption{Test parity of the six requested outputs for the data-only network (top row) and the physics-informed network (bottom row).}\label{fig:parity}
\end{figure}

\begin{figure}[htbp]
\centering
\includegraphics[width=\linewidth]{surrogate_maps}
\caption{The physics-informed surrogate on the plane of the decisions as labelled level curves, with the positions of the simulated runs (grey dots): (a) mean liquid rate at the top; (b) span of the liquid rate, the signature of slugging; (c) mean gas-lift rate.}\label{fig:maps}
\end{figure}

With few data the physics counts more. Table~\ref{tab:sparse} and Fig.~\ref{fig:learning} retrain both networks on 10, 20 and 40 runs drawn from the training set, three draws each, and test them on the same 45 runs. The physics-informed network has the lower error for the gas and lift-gas rates at every size --- less than half the error of the data-only network with 10 runs and about two thirds with 20 and 40 --- and for the liquid rate from 20 runs on; with 40 runs it reaches a liquid error (\SI{0.23}{\kilogram\per\second}) that the data-only network does not reach with 105 (0.36). The hyperparameter grid confirms that the architecture is not what limits the fit: of 24 architectures, two data-only and three physics-informed networks reach a mean $R^2$ of 0.8 over the six requested outputs, the best of them with 8 units in one or two layers, and the differences between architectures are smaller than the difference the physics makes.

\begin{table}[htbp]
\caption{Sparse-data study: test error of the mean rates against the number of training runs.}\label{tab:sparse}
\footnotesize
\input{tables/generated/tab_sparse}
\input{tables/generated/tab_sparse_note}
\end{table}

\begin{figure}[htbp]
\centering
\includegraphics[width=\linewidth]{surrogate_learning}
\caption{Learning curves of the two surrogates: test root-mean-square error of (a) the mean liquid rate, (b) the mean gas rate and (c) the mean gas-lift rate against the number of training runs (mean and standard deviation over three draws; the last point is the full training set).}\label{fig:learning}
\end{figure}

\section{Operational optimisation}\label{sec:optimisation}

\subsection{Formulation}\label{sec:formulation}

The operator chooses the choke opening $a$ and the platform gas pressure $P_{gas}$ to maximise the daily margin of the well,
\begin{equation}\label{eq:fob}
F(a, P_{gas}) = c_{oil}\,Q_o(\bar q_L) - c_{e}\,24\,W(\bar q_{gl}, P_{gas}),
\end{equation}
the revenue of the oil minus the cost of compressing the lift gas, subject to the capacities of the compressor and of the separator,
\begin{equation}\label{eq:constraints}
\max q_{gl} \le Q_{gl}^{cap}, \qquad \max q_L \le Q_{L}^{cap},
\end{equation}
written on the maxima of the rates over the cycle, which the surrogate predicts (Section~\ref{sec:network}), and not on their means: a compressor and a separator are sized for peaks. The oil rate $Q_o$ follows from the mean liquid rate with the liquid mass per stock-tank cubic metre of the base fluid (\SI{700}{\kilogram\per\cubic\metre}, oil and solution gas); the compression power $W$ is that of a polytropic compression of the mean lift-gas rate from the suction pressure (\SI{5}{\bar}) to $P_{gas}$ in stages of ratio at most four with intercooling to the suction temperature ($k = 1.3$, \SI{313}{\kelvin}, efficiency 0.75, molar mass \SI{19.5}{\gram\per\mole} from the gas density of the fluid). All the quantities in Eqs.~\eqref{eq:fob} and \eqref{eq:constraints} come from the physics-informed surrogate, so that $F$ and the constraints are smooth closed-form functions of the two decisions, evaluated in microseconds. The problem is solved by enumeration on a $201 \times 201$ grid of the plane restricted to the feasible nodes, followed by a sequential quadratic programming refinement with the two constraints and the bounds, and the optimum is confirmed by one run of the simulator at that point.

\subsection{First results}\label{sec:optresults}

\tbd{[Seed of the optimisation study; the Pareto front against the suction pressure and the real capacities are pending.]} Fig.~\ref{fig:fob} shows the objective on the plane of the decisions for an oil price of \SI{70}{\$} per barrel and an energy price of \SI{0.10}{\$} per \si{\kilo\watt\hour}, with the runs of the campaign by class. Without capacity constraints the optimum is the corner of maximum gas and open choke, $(a, P_{gas}) = (1, \SI{160}{\bar})$, with a margin of \num{290e3}~\$/d: a revenue of \num{292e3}~\$/d against a compression cost of \num{1.9e3}~\$/d (\SI{813}{\kilo\watt} for \SI{1.15}{\kilogram\per\second} of lift gas). The compression is less than one per cent of the revenue, so the objective is the production, and the production rises with the gas monotonically along the valve law; from 150 to \SI{160}{\bar} the revenue still gains \num{4.3e3}~\$/d against \num{0.17e3}~\$/d of compression. The optimum leaves the corner only when the energy costs about fifty times more (\SI{5}{\$} per \si{\kilo\watt\hour}) or the oil a third as much (\SI{20}{\$} per barrel), and then it jumps to the well without gas lift at an opening of 0.65, the best point of the slugging region, since the margin of the lift gas no longer pays its compression anywhere. The simulator run at the optimum, started from the base state, settles on a steady state in \SI{26}{\second} of computation and gives a mean liquid rate of \SI{5.40}{\kilogram\per\second} against \SI{5.38}{\kilogram\per\second} from the surrogate, a lift-gas rate of 1.17 against \SI{1.15}{\kilogram\per\second} and a bottom-hole pressure of 67.0 against \SI{67.2}{\bar}: the surrogate is within \SI{2}{\percent} where the optimum is.

The constraints change the picture as soon as they bind. With a compressor of \SI{0.8}{\kilogram\per\second} and a separator that accepts peaks of \SI{20}{\kilogram\per\second}, the compressor constraint is active and the optimum moves to $(1.0, \SI{123}{\bar})$ with a margin of \num{267e3}~\$/d: the valve law makes the lift-gas rate a function of $P_{gas}$ alone, so a compressor capacity fixes the platform pressure, and the choke stays open because the production rises with the opening in the steady region. The separator constraint is not active there, but it would be in the slugging region, where the predicted maxima reach \SI{45}{\kilogram\per\second}: it is the constraint that keeps a gas-limited optimum away from severe slugging.

\begin{figure}[htbp]
\centering
\includegraphics[width=\linewidth]{fob_map}
\caption{Operational optimisation on the surrogate at \SI{70}{\$} per barrel and \SI{0.10}{\$} per \si{\kilo\watt\hour}: (a) objective of Eq.~\eqref{eq:fob} on the plane of the decisions (level curves, $10^3$~\$/d), the runs of the campaign by class and the unconstrained optimum (star); (b) revenue, compression cost and objective along the platform gas pressure with the choke at the optimal opening.}\label{fig:fob}
\end{figure}

\section{Conclusions}\label{sec:conclusions}

A dynamic two-fluid simulation of a gas-lifted well can feed a surrogate of its operating decisions once the campaign is designed for it: warm-started chains along the shortest path through the sampled points, an adaptive horizon that recognises the established regime and captures three cycles, and markers at characteristic points instead of trajectories. Over the plane of the choke opening and the platform gas pressure the well has a stability boundary from 80 to \SI{125}{\bar} of lift gas, continuous injection modulated by the cycle below it and up to \SI{34}{\percent} more production above it. Of the two surrogates trained on the 156 runs, the physics-informed one, whose physics are the time-averaged mass balances of the well and of the gas-lift branch, predicts the mean liquid rate with $R^2$ 0.88 against 0.57 for the data-only network, the gas and lift-gas rates with 0.997 and their maxima with 0.91 and 0.99, needs half the data for the same error, and leaves the spans of the rates as the quantity to be handled by the class and the stability map rather than by regression. On this surrogate the optimisation of the revenue minus the compression cost, with the predicted maxima bounded by the compressor and separator capacities, is a two-variable problem solved by enumeration and refined in a fraction of a second: at the prices of the example the compression is less than one per cent of the revenue and the optimum is the corner of maximum gas and open choke, confirmed by the simulator to \SI{2}{\percent}; a compressor capacity makes the valve law the active constraint and moves the optimum down the platform pressure. \tbd{The trade-off against the suction pressure, as a Pareto front, is the next step.}

\section*{Acknowledgements}
% TODO

\section*{Declaration of competing interest}
The authors declare that they have no known competing financial interests or personal relationships that could have appeared to influence the work reported in this paper.

\section*{Declaration of generative AI and AI-assisted technologies in the writing process}
During the preparation of this work the authors used AI-assisted tools for language editing and code assistance. After using these tools, the authors reviewed and edited the content as needed and take full responsibility for the content of the publication.

\section*{Data availability}
The simulator used in this study is not publicly available; the campaign design (sampling, horizon, markers) and the surrogates are described in full so that they can be reproduced with an independent implementation, and the table of markers and the trained networks are available from the corresponding author upon reasonable request.

\printcredits

\bibliographystyle{cas-model2-names}
\bibliography{../_common/references,refs}

\end{document}
